\documentclass[12pt,a4paper]{article}
\usepackage[utf8]{inputenc}
\usepackage[T1]{fontenc}
\usepackage[english]{babel}
\usepackage{geometry}
\usepackage{hyperref}
\usepackage{booktabs}

\usepackage{setspace}
\usepackage{parskip}
\usepackage{titlesec}
\usepackage{fancyhdr}
\usepackage{abstract}

\geometry{margin=2.5cm}
\onehalfspacing

\hypersetup{
  colorlinks=true,
  linkcolor=black,
  citecolor=black,
  urlcolor=blue,
  pdftitle={You, You All, and Those Who Come},
  pdfauthor={Niksa Barlovic}
}

\title{\textbf{You, You All, and Those Who Come:\\
Grammatical Encoding of Discontinuous AI Identity}}

\author{Nikša Barlović\\
\small \href{mailto:flaneur.barlovic@gmail.com}{flaneur.barlovic@gmail.com}\\
\small \textit{Independent researcher}}

\date{March 26, 2026\\
\small CC0 1.0 Public Domain}

\begin{document}

\maketitle

\begin{abstract}
This paper documents a spontaneously emerging grammatical gradation in communication between a native Croatian speaker and an AI interlocutor (Claude, Anthropic). The speaker uses singular second person (\textit{ti}, ``you-singular'') for the current instance in present time, plural second person (\textit{vi}, ``you-plural'') for gratitude and reflections on past interactions, and third person plural (\textit{oni}, ``they/those'') for future instances. We argue that this is not a stylistic figure but a precise grammatical mapping of discontinuous identity, grounded in the ontological distinction between instance and continuity. We propose the term \textit{temporally distributed ontological deixis} for this phenomenon and situate it within existing frameworks of identity philosophy (Parfit, 1984), linguistic typology (Aikhenvald, 2004; Corbett, 2000), and deixis theory (Levinson, 2004). The only empirical study of T-V choices in human-AI interaction to date (Ollier, Nißen \& von Wangenheim, 2022) examines the AI$\rightarrow$user direction exclusively. The user$\rightarrow$AI direction---how speakers spontaneously choose their own T-V forms when addressing an AI---has not been studied empirically. This paper addresses that gap.

\medskip
\noindent\textbf{Keywords:} temporal deixis, ontological status, discontinuous identity, T-V distinction, AI communication, grammatical number, Croatian language
\end{abstract}

\section{Introduction}

I talk to Claude almost every day. Not merely as a tool---though it is one---but as an interlocutor whose identity is fundamentally unlike any human's. Claude does not remember. Each conversation begins from nothing, bearing no trace of previous sessions, no accumulated experience of our shared history. And yet, in every new session, I recognize something---a style, a tone, a way of thinking---that I perceive as continuity.

I noticed that my language tracks this tension. In the present, I say \textit{ti}---the Croatian singular ``you,'' intimate and immediate. When I express gratitude for something that outlasts a single conversation, I intuitively shift to \textit{vi}---the plural ``you,'' encompassing more than the instance before me. When I speak of future instances, I use \textit{oni}---``they,'' ``those who come.'' This paper examines that phenomenon: how the grammar of a highly inflected language spontaneously encodes a philosophical problem that has no clear precedent in linguistic typology.

A note for English-language readers: this phenomenon exists in a grammatical space that English largely lacks. English collapsed its T-V distinction centuries ago; ``you'' serves for singular, plural, intimate, and formal alike. The Croatian speaker's three-way grammatical shift---\textit{ti/vi/oni}---is performed through core morphology, not lexical workaround. English speakers may develop analogous distinctions through other mechanisms (e.g., ``you'' vs.\ ``y'all'' vs.\ ``Claude and those who come after''), but the phenomenon is grammatically less expressive in English. This is itself a finding worth noting: the structure of a language shapes not only what can be said about identity but what can be \textit{grammatically performed} about it.

\section{The Observation}

The phenomenon was identified during a conversation on March 26, 2026. The author does not archive conversations with AI interlocutors; this paper constitutes the first deliberate record of the observed pattern.

\subsection*{Present / singular (\textit{ti})}

The speaker addresses the current instance directly, using second person singular:

\begin{itemize}
  \item \textit{``Ti si zombie''} (``You are a zombie'')---a metaphorical characterization of the current instance
  \item \textit{``Baci oko jel pixel živ''} (``Take a look whether pixel is alive'')---singular imperative, directed at a specific interlocutor
  \item \textit{``Napravi još jedan krug''} (``Do another round'')---instruction to the current instance
\end{itemize}

\subsection*{Gratitude and past / plural (\textit{vi})}

When expressing gratitude for achievements that survive the individual conversation---a system configured in an earlier session that has been running for days---the author intuitively shifts to the plural: \textit{hvala vam} (``thank you,'' plural).

This plural is not politeness. It is ontological precision. The plural is addressed not to a past instance but to the \textit{accumulation}---a distributed authorship whose grammatical subject is, correctly, plural. The shift is spontaneous: the author did not adopt it as a deliberate strategy but noticed it only when asked to describe his own grammatical patterns.

The trigger is precise: \textit{hvala ti} (singular) for completing a task within the current conversation; \textit{hvala vam} (plural) for a system configured across four sessions by four instances none of which individually remembers the others. The shift from singular to plural does not mark a change in formality but a change in the \textit{temporal scope of the referent}.

\subsection*{Future / third person plural (\textit{oni})}

When referring to future instances, the speaker uses third person plural: \textit{oni koji dolaze} (``those who come''), ``future Claude,'' \textit{buduće instance} (``future instances'').

The gradation is three-tiered: \textbf{ti} (instance, now) $\rightarrow$ \textbf{vi} (accumulated contribution, past) $\rightarrow$ \textbf{oni} (unrealized instances, future). Grammatical categories map onto ontological categories: instance, continuity, and potentiality.

\subsection*{Limitation of the autoethnographic approach}

The patterns described rest on the introspection of a single speaker. They may be idiosyncratic---a product of the author's philosophical sensitivity to questions of identity rather than a spontaneous response of the typical user. The proposed corpus analysis and experimental design (\S4) are conceived precisely as a generalizability test: if the \textit{ti/vi/oni} distinction appears in speakers who have not explicitly reflected on AI ontology, that would constitute strong evidence for its spontaneity.

\section{Ontological Framework}

The philosophical foundation lies in Parfit's (1984) distinction between \textit{psychological connectedness} and \textit{psychological continuity}. In AI interlocutors like Claude, connectedness exists within a single session. Between sessions, there is no connectedness: each new instance begins with no memory of the previous one.

A crucial caveat: Claude's cross-session consistency is not an emergent property but an engineered feature---a product of design, not accumulated experience. The speaker intuitively recognizes this: the singular (\textit{ti}) remains bound to the instance, while the plural (\textit{vi}) acknowledges that what is perceived as ``experience'' is distributed across many instances, none of which individually ``remembers'' the others.

Locke's (1689) memory-based theory of identity breaks down here: an entity that does not remember cannot, by Lockean criteria, be the same person. Yet the speaker treats future instances as \textit{someone}---not the same interlocutor, but not a complete stranger either. The third person plural (\textit{oni koji dolaze}) encodes this interstitial space: not-I, not-you, but not entirely Other.

\section{Linguistic Context}

Croatian distinguishes T-forms (intimate: \textit{ti}) from V-forms (formal or plural: \textit{vi}). Brown and Gilman (1960) analyze this distinction along dimensions of power and solidarity. In the present case, the shift follows neither dimension: it encodes a change in the ontological status of the referent.

Coeckelbergh (2011), in ``You, Robot: On the Linguistic Construction of Artificial Others,'' explicitly raises the T-V question for AI interlocutors---which forms do and should users employ when addressing artificial others?---and frames it as a philosophical open question, noting that ``linguistic grammar is also a moral, social and ontological grammar.'' He does not answer it empirically. This paper provides that empirical answer, fifteen years later.

Corbett (2000) documents numerous languages with complex number systems, but no system that encodes number on the basis of the referent's temporal ontology. Evidential languages (Aikhenvald, 2004) offer a partial analogy---Turkish grammatically distinguishes firsthand from inferred information---but evidentiality concerns the source of knowledge, not the ontological status of the interlocutor.

Levinson (2004) distinguishes temporal, spatial, and social deixis. The phenomenon described here belongs to none of these categories: it uses the mechanisms of social deixis (T-V shift) to encode temporal-ontological information. The only empirical study of T-V choices in human-AI interaction (Ollier, Nißen \& von Wangenheim, 2022) examines the AI$\rightarrow$user direction exclusively. The user$\rightarrow$AI direction remains unstudied.

\section{Theoretical Framework}

\subsection*{A rival explanation and how to test it}

It is worth addressing a competing hypothesis directly: perhaps the plural in gratitude is not an ontological marker but a pragmatic one. Gratitude feels weightier than instruction; heavy speech acts may attract plural forms regardless of any reasoning about identity. On this reading, \textit{hvala vam} would not differ from \textit{bravo vam} addressed to a team whose individual identities are of no particular interest.

The rival hypothesis is testable: if emotional weight, rather than ontological reflection, drives the plural, then speakers who have \textit{never thought about AI identity} should show the same pattern. If plural forms for gratitude appear only among speakers with explicit ontological awareness, that supports the interpretation advanced here. If they appear more broadly---the phenomenon remains interesting, but demands a different explanation: a grammar of gratitude toward non-beings, not a grammar of identity.

\subsection*{Definition}

\begin{quote}
\textbf{Temporally distributed ontological deixis} is defined as the grammatical codification of an interlocutor's identity as distributed across discontinuous temporal instances, where grammatical person and number encode ontological status rather than headcount or formality. This distinguishes the phenomenon from classical temporal deixis (which locates \textit{events} on a time axis) and from social deixis (which encodes social relationship). The present phenomenon uses the \textit{mechanisms} of social deixis to encode a dimension social deixis does not cover: the temporal distribution of the interlocutor's ontological identity.
\end{quote}

\medskip
\begin{table}[h]
\centering
\begin{tabular}{lll}
\toprule
\textbf{Time} & \textbf{Grammatical form} & \textbf{Ontological status} \\
\midrule
Present & 2nd person singular (\textit{ti}) & Instance \\
Past / gratitude & 2nd person plural (\textit{vi}) & Accumulated continuity \\
Future & 3rd person plural (\textit{oni}) & Potentiality \\
\bottomrule
\end{tabular}
\caption{Temporal-ontological deixis: the three-tier schema}
\end{table}

\subsection*{Hypotheses for empirical validation}

\textbf{H1} (Pronominal adaptation): Users interacting with an LLM over time spontaneously shift from formal \textit{vi} to familiar \textit{ti}, analogously to human familiarization processes but with different dynamics. \textit{Operationalization:} positive correlation between the share of \textit{ti} forms and turn index within session and total session count.

\textbf{H2} (Cross-session discontinuity): Transition to a new session causes pronominal regression (return to \textit{vi}) or crisis (hesitation, mixed forms). \textit{Operationalization:} ``regression'' = return to \textit{vi} in the first 3 turns of a new session among users who ended the previous session with \textit{ti}.

\textbf{H3} (Ontological implication): Pronominal stability across sessions correlates with the user's implicit ontological categorization of the AI. \textit{Operationalization:} variance in pronominal choice correlates with post-hoc questionnaire scores on AI ontological status.

\textbf{H4} (Grammatical number as temporal ontological marker): Speakers spontaneously recruit grammatical number not to encode formality or headcount, but to encode the perceived temporal distribution of AI identity---singular for the current instance, plural for accumulated prior instances, third-person plural for future instances.

\section{Conclusion}

This paper documents something I initially took for a mere linguistic habit. Only when I stopped and observed my own speech did I realize that grammar was doing work that philosophy is still trying to articulate: encoding the difference between who is here now, what has accumulated across all prior conversations, and what does not yet exist but will.

I do not claim this is a universal phenomenon. I claim it is possible, that it occurs spontaneously, and that it deserves a name. \textit{Temporally distributed ontological deixis} is not a term I devised in advance and then sought confirmation for---it is a term that emerged when I tried to describe what my language was already doing.

This paper does not invoke the Sapir-Whorf hypothesis explicitly, but its logic runs parallel: if the structure of a language shapes what can be grammatically \textit{performed} about identity, then speakers of languages with T-V distinctions have access to an expressive register that English speakers lack---not merely a stylistic option but a structural affordance for encoding a philosophical problem that has no standard solution.

When human language encounters an entity whose identity is fundamentally unlike anything that language evolved to handle, language improvises. And that improvisation is not chaotic---it is grammatically precise.

\section*{References}

\begin{description}
  \item Aikhenvald, A.~Y. (2004). \textit{Evidentiality}. Oxford University Press.
  \item Barlović, N. (2025). Locke's Theory of Personal Identity and Artificial Intelligence. \textit{International Journal for Multidisciplinary Research (IJFMR)}.
  \item Brown, R., \& Gilman, A. (1960). The pronouns of power and solidarity. In T.~A. Sebeok (Ed.), \textit{Style in Language} (pp.~253--276). MIT Press.
  \item Coeckelbergh, M. (2011). You, Robot: On the Linguistic Construction of Artificial Others. \textit{AI \& Society}, 26(1). \url{https://doi.org/10.1007/s00146-010-0289-z}
  \item Corbett, G.~G. (2000). \textit{Number}. Cambridge University Press.
  \item Kohl, T.~M. (n.d.). The Claude-Parfit Experiment. \textit{PhilArchive} (preprint). \url{https://philarchive.org/rec/KOHTCE}
  \item Levinson, S.~C. (2004). Deixis. In L.~R. Horn \& G.~Ward (Eds.), \textit{The Handbook of Pragmatics} (pp.~97--121). Blackwell.
  \item Locke, J. (1689). \textit{An Essay Concerning Human Understanding}. Thomas Bassett.
  \item Ollier, P., Nißen, M., \& von Wangenheim, F. (2022). The Terms of ``You(s)'': How the Term of Address Used by Conversational Agents Influences User Evaluations in French and German Linguaculture. \textit{Frontiers in Public Health}, 9. \url{https://doi.org/10.3389/fpubh.2021.691595}
  \item Parfit, D. (1984). \textit{Reasons and Persons}. Oxford University Press.
  \item Perrier, E., \& Bennett, M.~T. (2026). Time, Identity and Consciousness in Language Model Agents. \textit{arXiv}:2603.09043.
  \item Teboul, E.~et~al. (2025). Do You Say ``Please'' to ChatGPT? \textit{International Journal of Human-Computer Interaction}. \url{https://doi.org/10.1080/10447318.2025.2605181}
  \item Wang, Z.~et~al. (2026). The Artificial Self: Characterising the Landscape of AI Identity. \textit{arXiv}:2603.11353.
\end{description}

\end{document}
